\section{Proof Sketch}

The argument begins with a coefficient-envelope calculation on the actual
unit box.  Lemma~\ref{lem:unit-box-envelope} bounds polynomial values and
partials by their monomial coefficient $\ell_1$ budgets, and
Proposition~\ref{prop:coordinate-envelope} applies the common-chain rule to
obtain $D_*=\Delta B_Q(1+qB_P)$.  The literal anchor then supplies a global
denominator.  Lemma~\ref{lem:normalized-curve} computes the Euclidean
projector derivative, and Proposition~\ref{prop:projective-speed} transfers it
through $x=(\theta-c)/h$ with the exact factor $h^{-1}$.

For the central event, Lemma~\ref{lem:incidence-jacobians} computes the two
tangential Jacobians of the incidence hypersurface.  The area and coarea
formulas \citep{federer1969gmt} turn projection multiplicity into the
coefficient-volume inequality in Proposition~\ref{prop:central-volume}; the
null classes and boundary regimes are handled in
Lemma~\ref{lem:central-null-classes}.  Ball's sharp central cube-section
theorem \citep{ball1986cube} is scaled exactly in
Lemma~\ref{lem:scaled-cube-section}.  Proposition~\ref{prop:central-sweep}
then uses the single full joint-density cap, and
Proposition~\ref{prop:central-rate} performs the visible speed substitution
and the interval-then-law supremum closure.

The affine argument solves the root equation in a selected nonzero
coefficient coordinate.  Lemmas~\ref{lem:finite-chart-legality} and
\ref{lem:affine-chart-jacobian} establish finite-level Lipschitz legality and
the exact determinant $|\partial_\theta T_j|$.
Proposition~\ref{prop:finite-chart-area} applies the equal-dimensional area
formula, while Lemma~\ref{lem:affine-exhaustion} proves that every root enters
a least finite pivot level.  Monotone convergence yields the extended-real
bound in Proposition~\ref{prop:affine-bound}, without transversality or a
uniform pivot margin.

For monic polynomials, Proposition~\ref{prop:monic-presentation} constructs
the exact normalized presentation while keeping the leading coefficient in
the deterministic offset.  Lemma~\ref{lem:monic-pivot-partition} supplies the
prescribed two cells, and Lemmas~\ref{lem:monic-low-chart} and
\ref{lem:monic-high-chart} compute their velocities.  The general affine
theorem is transferred without loss in
Proposition~\ref{prop:monic-affine-transfer};
Lemma~\ref{lem:monic-two-cell-ledger} keeps the sharper high-chart constant
until the final coefficient-one comparison, while
Proposition~\ref{prop:monic-linear-branch} handles the zero-dimensional
$d=1$ case.

Finally, Proposition~\ref{prop:counter-presentation} computes the exact
projective speed of the two-feature example,
Lemma~\ref{lem:counter-wedges} identifies both closed coefficient wedges and
their area, and Proposition~\ref{prop:counter-scale} compares the resulting
$1/(4\delta)$ certificate with the distinct projective and raw-presentation
upper scales.  Proposition~\ref{prop:full-conjunction} places these exact
interfaces under their respective numbered assumptions, after which the main
theorem is their direct conjunction.
